\documentclass[11pt]{amsart}

\usepackage[T1]{fontenc}
\usepackage{lmodern}
\usepackage{microtype}
\usepackage{amsmath,amssymb,amsthm}
\usepackage[hidelinks]{hyperref}

\hypersetup{
  pdftitle={Counterexamples to the Xu--He--Lu Conjecture on Partite Cycle Saturation}
}

\newtheorem{theorem}{Theorem}
\newtheorem{lemma}[theorem]{Lemma}

\newcommand{\sat}{\operatorname{sat}}

\title[The Xu--He--Lu partite saturation conjecture]
{Counterexamples to the Xu--He--Lu Conjecture on Partite Cycle Saturation}

\date{}

\subjclass[2020]{05C35}
\keywords{graph saturation, partite saturation, bipartite graph, even cycle}

\begin{document}

\begin{abstract}
Xu, He, and Lu conjectured that, for $\ell\ge 4$ and
$n_1,n_2\ge \ell+2$,
\[
\sat(K_{n_1,n_2},C_{2\ell})
 =n_1+n_2+\ell^2-3\ell+1.
\]
For every $\ell\ge 4$, we construct a $C_{2\ell}$-saturated spanning
subgraph of $K_{2\ell-2,2\ell-2}$ with $\ell^2-1$ edges. The
conjectured value at these parameters is $\ell^2+\ell-3$, so the
construction has $\ell-2$ fewer edges than the conjectured value and
refutes the conjecture. The smallest example is a $15$-edge
$C_8$-saturated subgraph of $K_{6,6}$, where the conjecture predicts
$17$ edges.
\end{abstract}

\maketitle

\section{Result}\label{sec:result}

For graphs $F$ and $G$, a spanning subgraph $H\subseteq G$ is
\emph{$F$-saturated relative to $G$} if $H$ is $F$-free but $H+e$
contains a copy of $F$ for every $e\in E(G)\setminus E(H)$. The
minimum number of edges in such a graph is denoted by $\sat(G,F)$.

Xu, He, and Lu~\cite[Theorem~1.1]{XHL} proved that
\begin{equation}\label{eq:XHLupper}
\sat(K_{n_1,n_2},C_{2\ell})
 \le n_1+n_2+\ell^2-3\ell+1
\end{equation}
for $\ell\ge 3$ and $n_1,n_2\ge \ell+2$, and conjectured that this
upper bound is the exact value. Verbatim,
\cite[Conjecture~6.2]{XHL} states: ``For $\ell\ge 4$ and
$n_1,n_2\ge \ell+2$,
$\sat(K_{n_1,n_2},C_{2\ell})=n_1+n_2+\ell^2-3\ell+1$.'' For reference
below we display the asserted equality,
\begin{equation}\label{eq:XHL}
\sat(K_{n_1,n_2},C_{2\ell})
 =n_1+n_2+\ell^2-3\ell+1
\qquad(\ell\ge 4,\ n_1,n_2\ge \ell+2).
\end{equation}

\begin{theorem}\label{thm:main}
For every integer $\ell\ge 4$, there is a $C_{2\ell}$-saturated
spanning subgraph $H_\ell\subseteq K_{2\ell-2,2\ell-2}$ with
\[
e(H_\ell)=\ell^2-1.
\]
Consequently,
\[
\sat(K_{2\ell-2,2\ell-2},C_{2\ell})\le \ell^2-1,
\]
and \cite[Conjecture~6.2]{XHL} is false for every $\ell\ge 4$.
\end{theorem}

Indeed, $2\ell-2\ge \ell+2$, while the right-hand side of
\eqref{eq:XHL} at $n_1=n_2=2\ell-2$ is
$\ell^2+\ell-3$, exceeding $\ell^2-1$ by $\ell-2$. Thus
Theorem~\ref{thm:main} also improves the upper
bound~\eqref{eq:XHLupper} by $\ell-2$ edges at $n_1=n_2=2\ell-2$,
which for $\ell=4$ is the smallest admissible pair
$n_1=n_2=\ell+2=6$. It refutes~\eqref{eq:XHL} only at these balanced
parameters; whether the equality in~\eqref{eq:XHL} holds when $n_1$
and $n_2$ are large relative to $\ell$ remains open.

\section{Construction and proof}

All path lengths below count edges. Fix $\ell\ge 4$ and put
$r=\ell-1$. Let
\[
A_0=\{a_0,a_1,\ldots,a_{r-1}\},\qquad
B_0=\{b_0,b_1,\ldots,b_{r-1}\},
\]
and
\[
X=\{x_1,\ldots,x_r\},\qquad
Y=\{y_1,\ldots,y_r\}.
\]
Take $A=A_0\cup X$ and $B=B_0\cup Y$ as the two partite classes.
Let
\[
R=K_{r,r}-a_0b_0
\]
on $A_0\cup B_0$, and add the edges of the path
\begin{equation}\label{eq:Q}
Q\colon a_0-y_1-x_1-y_2-x_2-\cdots-y_r-x_r-b_0.
\end{equation}
Call the resulting graph $H_\ell$. Its two parts have size $2r$, and
\begin{equation}\label{eq:edgecount}
e(H_\ell)=(r^2-1)+(2r+1)=r^2+2r=\ell^2-1.
\end{equation}
For vertices $u,v$ on $Q$, write $Q[u,v]$ for the unique $u$--$v$
subpath of $Q$ and $|Q[u,v]|$ for its length.

\begin{lemma}\label{lem:core}
The graph $R$ contains
\begin{enumerate}
\item an $a_0$--$b_0$ path of every odd length $2s+1$ with
      $1\le s\le r-1$;
\item for each $a\in A_0\setminus\{a_0\}$, an $a$--$a_0$ path of
      every even length $2s$ with $1\le s\le r-1$.
\end{enumerate}
The symmetric statement in $B_0$ also holds.
\end{lemma}

\begin{proof}
For the first assertion, choose distinct
$\alpha_1,\ldots,\alpha_s\in A_0\setminus\{a_0\}$ and distinct
$\beta_1,\ldots,\beta_s\in B_0\setminus\{b_0\}$. Then
\[
a_0-\beta_1-\alpha_1-\beta_2-\alpha_2-\cdots-
\beta_s-\alpha_s-b_0
\]
is a path of length $2s+1$.

For the second, choose distinct
$\beta_1,\ldots,\beta_s\in B_0\setminus\{b_0\}$ and distinct
$\alpha_1,\ldots,\alpha_{s-1}\in
A_0\setminus\{a_0,a\}$. The sequence of $\alpha_i$ is empty when
$s=1$, and
\[
a-\beta_1-\alpha_1-\beta_2-\cdots-
\alpha_{s-1}-\beta_s-a_0
\]
is a path of length $2s$. The assertion in $B_0$ follows by symmetry.
\end{proof}

\begin{proof}[Proof of Theorem~\ref{thm:main}]
We first show that $H_\ell$ is $C_{2\ell}$-free. A cycle contained
in $R$ has length at most $2r=2\ell-2$. Every internal vertex of $Q$
has degree two and no neighbor outside $Q$. Hence a cycle meeting an
internal vertex of $Q$ must traverse all of $Q$ and then return from
$b_0$ to $a_0$ through $R$. Since $a_0b_0\notin E(R)$, such a core
path has length at least three. The cycle therefore has length at
least
\[
(2r+1)+3=2r+4=2\ell+2.
\]
Thus $H_\ell$ contains no $C_{2\ell}$.

It remains to show that the endpoints of every edge in
$E(K_{2r,2r})\setminus E(H_\ell)$ are joined in $H_\ell$ by a simple
path of length $2r+1=2\ell-1$. Adding the missing edge then closes
this path to a $C_{2\ell}$. Along $Q$ we have
\begin{align*}
|Q[a_0,y_j]|&=2j-1, & |Q[y_j,b_0]|&=2r-2j+2,\\
|Q[a_0,x_i]|&=2i,   & |Q[x_i,b_0]|&=2r-2i+1.
\end{align*}

The missing edge $a_0b_0$ is witnessed by $Q$ itself. Next consider
a missing edge $ay_j$ with $a\in A_0$.
If $a\ne a_0$ and $j=1$, concatenate the edge $ab_0$ with
$Q[b_0,y_1]$; its length is $1+2r$.
If $a\ne a_0$ and $j\ge2$, Lemma~\ref{lem:core} gives an
$a$--$a_0$ path in $R$ of length $2(r-j+1)$, since
$1\le r-j+1\le r-1$. Appending $Q[a_0,y_j]$ gives total length
\[
2(r-j+1)+(2j-1)=2r+1.
\]
If $a=a_0$ and $j\ge2$, take an $a_0$--$b_0$ path in $R$ of length
$2j-1$ and append $Q[b_0,y_j]$; the total length is
\[
(2j-1)+(2r-2j+2)=2r+1.
\]
Each concatenation is simple because its core portion meets $Q$ only
at the indicated endpoint or endpoints.

The map determined by
\[
a_0\leftrightarrow b_0,\qquad
a_k\leftrightarrow b_k\ (1\le k\le r-1),\qquad
x_i\leftrightarrow y_{r-i+1}\ (1\le i\le r)
\]
preserves $R$ and reverses $Q$, and hence is an automorphism of
$H_\ell$. It maps the missing edges between $A_0$ and $Y$ to those
between $X$ and $B_0$, so the latter are handled as well.

Finally, consider a missing edge $x_i y_j$. The edges with $j=i$,
and with $j=i+1$ when $i<r$, already lie on $Q$.
If $j\le i-1$, concatenate $Q[x_i,b_0]$, a $b_0$--$a_0$ path in $R$
of length $2(i-j)+1$, and $Q[a_0,y_j]$. Lemma~\ref{lem:core}
applies because $1\le i-j\le r-1$, and the total length is
\[
(2r-2i+1)+(2i-2j+1)+(2j-1)=2r+1.
\]
If $j\ge i+2$, concatenate $Q[x_i,a_0]$, an $a_0$--$b_0$ path in $R$
of length $2(j-i)-1$, and $Q[b_0,y_j]$. Here
$1\le j-i-1\le r-2$, and the total length is
\[
2i+(2j-2i-1)+(2r-2j+2)=2r+1.
\]
In both cases the two subpaths of $Q$ are vertex-disjoint, while the
core path meets $Q$ only at $a_0$ and $b_0$; hence the concatenation is
simple. These cases exhaust all missing edges of the host graph, so
$H_\ell$ is $C_{2\ell}$-saturated. Equation~\eqref{eq:edgecount}
completes the proof.
\end{proof}

For $\ell=4$, the construction is $K_{3,3}-a_0b_0$ together with a
seven-edge $a_0$--$b_0$ path. It is a $15$-edge $C_8$-saturated
spanning subgraph of $K_{6,6}$, whereas~\eqref{eq:XHL} gives $17$.

\begin{thebibliography}{9}

\bibitem{XHL}
Y.~Xu, Z.~He, and M.~Lu,
\emph{Partite saturation number of cycles},
Discrete Math. \textbf{349} (2026), no.~3, Paper No.~114802.
\href{https://doi.org/10.1016/j.disc.2025.114802}
{doi:10.1016/j.disc.2025.114802};
\href{https://arxiv.org/abs/2410.11194}{arXiv:2410.11194}.
The theorem and conjecture numbers cited here, and the quotation in
Section~\ref{sec:result}, follow version~2 of the e-print.

\end{thebibliography}

\end{document}